% ID: SEL-D0075-OB-E03
% Tanda S075.
\textbf{Parte I.} Denotamos por $S$ tener estudios superiores y por $E$ tener empleo.

\textbf{a)} Separamos los dos grupos de población. Por la probabilidad total,
\[
 P(E)=P(E\mid S)P(S)+P(E\mid S^c)P(S^c)
 =0{,}95\cdot0{,}30+0{,}60\cdot0{,}70=\boxed{0{,}705}.
\]
\textbf{b)} Por Bayes, $P(S\mid E)=0{,}285/0{,}705=\boxed{19/47}\approx0{,}4043$.

\textbf{Parte II. a)} Tomamos muestreo aleatorio simple sin reemplazo. Como el orden no afecta a la media, hay $\binom42=6$ pares equiprobables:
\[
 \begin{array}{c|cccccc}
 \text{Muestra}&\{1,2\}&\{1,3\}&\{1,4\}&\{2,3\}&\{2,4\}&\{3,4\}\\
 \bar X&1{,}5&2&2{,}5&2{,}5&3&3{,}5
 \end{array}
\]
\textbf{b)} La esperanza de las medias es $2{,}5$. La varianza de su distribución es el promedio ponderado de las desviaciones al cuadrado:
\[
 \operatorname{Var}(\bar X)=\frac{1+0{,}25+0+0+0{,}25+1}{6}
 =\boxed{\frac5{12}}.
\]

% AUD001: alternativa para hipótesis no explícita.

\textit{Hipótesis de muestreo.} El enunciado no explicita la reposición.
Si se distingue el orden sin reemplazo, las doce parejas equiprobables
producen la misma distribución de medias que los seis subconjuntos.

Con reemplazo, en cambio, hay dieciséis pares ordenados equiprobables.
En la tabla, la fila representa la primera extracción y la columna
la segunda. Cada celda da la media del par:
\[
\begin{array}{c|cccc}
i\backslash j&1&2&3&4\\\hline
1&1&1{,}5&2&2{,}5\\
2&1{,}5&2&2{,}5&3\\
3&2&2{,}5&3&3{,}5\\
4&2{,}5&3&3{,}5&4
\end{array}
\]
La esperanza sigue siendo $2{,}5$. La varianza poblacional es
$[(1-2{,}5)^2+(2-2{,}5)^2+(3-2{,}5)^2+(4-2{,}5)^2]/4=5/4$.
Por independencia de las dos extracciones,
$\operatorname{Var}(\bar X)=(5/4)/2=\boxed{5/8}$.
La elección entre $5/12$ y $5/8$ depende de la hipótesis de reposición.
